\chapter{From Minimum Wage to Real Estate Millionaire}

This chapter follows Alex Kaufman's interview with School of Hard Knocks in the series \emph{10 Questions with a Millionaire}, curated by LazyingArt LLC through Video2Book. The subject is not a blackboard theory of markets, but it has a quantitative spine: income became ownership, ownership became equity, equity produced optionality, and risk was treated not as a slogan but as a tradeoff between possible loss and forfeited upside. We will keep the story in the order in which it is told: first the personal bottom, then the business definition, then the first deal, then partnership, risk, and the larger purpose of the enterprise.

\section{The Reversal Before The Resume}

The interview begins before the formal introduction, with the most severe fact placed first. Alex says that from about age eleven to twenty-one he was under the influence of alcohol and drugs, and that by twenty-one he was a heroin addict. Only after that does the host reset the scene: this is episode five of \emph{10 Questions with a Millionaire}; Alex Kaufman is introduced as a real estate brokerage owner and real estate investor.

That ordering matters. If we wrote the story only as a business summary, we would miss the first transformation. The starting state is not ``young entrepreneur discovers real estate.'' It is closer to

\begin{equation}
\text{personal disorder}
\quad \longrightarrow \quad
\text{sobriety}
\quad \longrightarrow \quad
\text{legal enterprise}.
\end{equation}

Alex's first explicit business facts come quickly. He says he got started in real estate about five years earlier with his partner, Matt Teifke. They began by investing together, then partnered on the brokerage side, scaled in Texas, and described the goal as a global real estate brokerage for entrepreneurs. At the time of the interview, he states that they own approximately

\begin{equation}
U_{\text{owned}} \approx 100
\end{equation}

units together. This is the first scale marker: the chapter begins with personal ruin, but the present-tense credential is ownership.

\section{The Brokerage As A Multi-Lane Skill Stack}

The host then asks what Alex means by ``entrepreneur.'' Alex answers by contrasting two models of real estate training. In the ordinary brokerage path, a new licensee may be taught to be one thing: a listing agent, a buyer's agent, a residential agent, or a commercial agent. Alex's criticism is not that any one of these skills is useless. It is that a single skill can become a trap when the market changes.

We can express the distinction as a simple model of channels. Let

\begin{equation}
\mathcal{C}
=
\{\text{listing},\text{buyer agency},\text{residential},\text{commercial},\text{investing},\text{wholesaling}\}.
\end{equation}

A one-lane agent operates on a small subset of this set, sometimes only one element. The TRE entrepreneur, as Alex describes it, is trained to make money ``in and on'' real estate across more of the set.

\begin{figure}
\centering
\begin{tikzpicture}[
  box/.style={draw, rounded corners, align=center, text width=0.76\linewidth, minimum height=1.1cm},
  arrow/.style={->, thick}
]
\node[box] (traditional) {Traditional brokerage training};
\node[box, below=0.55cm of traditional] (lane) {One narrow lane\\listing, buyer agency, residential, or commercial};
\node[box, below=0.55cm of lane] (shift) {Market shift};
\node[box, below=0.55cm of shift] (limit) {Limited earning capacity\\if the one lane weakens};

\draw[arrow] (traditional) -- (lane);
\draw[arrow] (lane) -- (shift);
\draw[arrow] (shift) -- (limit);
\end{tikzpicture}
\caption{The narrow-path risk described in the interview: specialization becomes fragile when the market shifts.}
\end{figure}

\subsection{Question \& Answer}

\paragraph{Question.}
Why is being taught only one real estate lane a business risk?

\paragraph{Answer.}
Because a market shift changes which channels are available. If an agent only knows one way to earn, then the agent's income is tied to the condition of that one path. Alex's answer is to define the real estate entrepreneur as someone with a wider operating set:

\begin{equation}
\text{earning capacity}
\approx
f(\text{number of usable channels},\ \text{market condition},\ \text{execution}).
\end{equation}

This is not a formal law. It is the practical mechanism he is describing: a broader skill stack gives more ways to survive and earn when one activity stops working.

\section{Sobriety, Minimum Wage, And The Legal Money Path}

After defining the business model, the host asks for Alex's career path. The answer returns to the biographical opening. Alex got sober at twenty-one. He worked at Firehouse Subs for

\begin{equation}
w_{\text{Firehouse}} = \$8.50/\text{hour}.
\end{equation}

He then washed dishes at a country club, worked in the kitchen, and worked for an oil and gas company. The point is not merely that the jobs were humble. The point is that the jobs did not match the direction he believed he needed to go. He says he knew he did not want a job or to follow a conventional career path, and that he believed he could and would make money for himself.

But after sobriety, the question changed. The old illegal route had to be replaced. The new question was: how can one make money legally, with discipline, and without returning to the destructive channel?

That question leads to Matt Teifke. Alex says he had known Matt since he was about eight years old, and Matt showed him the path in real estate. The first transition into business is therefore not abstract education. It is relationship plus trust.

\section{Matt, Trust, And The First Deal}

When the host asks how Alex moved from starting in real estate to becoming involved with a brokerage, Alex clarifies that Matt is the licensed broker, but they own the brokerage together. The brokerage is formed as a counter-model to the narrow paths described earlier. They wanted a place where people could get more out of real estate for themselves.

At that point, Alex says they had already been buying real estate together and owned about fifteen units. The chronology is important:

\begin{enumerate}
\item They invest together first.
\item They see the limits of conventional brokerage paths.
\item Matt has the broker license.
\item They build TRE as a brokerage for real estate entrepreneurs.
\end{enumerate}

The transcript is garbled in the lead-in to the first-deal story, so we should move carefully. The clear elements are these: Alex had never bought real estate before; the amount stated is about \$28,000; Matt told him to trust the deal; Alex placed trust in Matt's judgment; the property later became, in Alex's spoken estimate, worth close to \$400,000; they had pulled about \$150,000 out of it; they still owned it; and it still produced monthly rent cash flow.

\subsection{Question \& Answer}

\paragraph{Question.}
How could the first deal work when Alex had never bought real estate before?

\paragraph{Answer.}
In the story as told, the missing technical experience was partly replaced by trusted partnership. Matt supplied enough judgment for Alex to act. That does not make trust a substitute for underwriting in general; it explains the local mechanism of this first step. The practical formula is

\begin{equation}
\text{first action}
=
\text{trust}
+
\text{available capital}
+
\text{partner's domain judgment}.
\end{equation}

The chapter should not turn that into a universal rule. Here it is an anecdote: the first acquisition depended on Matt's conviction and Alex's willingness to act despite having no prior purchase experience.

\section{The First Deal As A Wealth Mechanism}

Now we can isolate the strongest quantitative payload of the interview. Let \(C_0\) denote the stated first-deal cash figure. We do not know from the transcript whether it was a down payment, total cash needed, a contribution, or some other capital figure, so we name it cautiously:

\begin{equation}
C_0 \approx \$28{,}000.
\end{equation}

Let \(V_{\text{later}}\) denote Alex's spoken estimate of the later property value:

\begin{equation}
V_{\text{later}} \approx \$400{,}000.
\end{equation}

Let \(E_{\text{out}}\) denote the amount he says they pulled out while still owning it:

\begin{equation}
E_{\text{out}} \approx \$150{,}000.
\end{equation}

A rough scale comparison is useful, provided we do not overstate it:

\begin{align}
\frac{V_{\text{later}}}{C_0}
&\approx
\frac{400{,}000}{28{,}000} \\
&\approx 14.3.
\end{align}

This is not a complete return calculation. It does not include debt, purchase price, taxes, repairs, interest, closing costs, time, or operating expenses. It is only a ratio between two spoken figures: the stated initial cash figure and the later estimated value.

\subsection{A Worked Mechanism}

The important part of the story is not just that the value rose. It is that the ownership was not given up. We can write the mechanism in steps.

\begin{enumerate}
\item Initial capital is committed:
\[
C_0 \approx \$28{,}000.
\]

\item The property is acquired and retained.

\item Alex later estimates its value near
\[
V_{\text{later}} \approx \$400{,}000.
\]

\item Some value is extracted:
\[
E_{\text{out}} \approx \$150{,}000.
\]

\item Ownership continues, and the property still produces rent cash flow.
\end{enumerate}

So the financial pattern is better represented as

\begin{equation}
\text{wealth effect}
\approx
\text{retained ownership}
+
\text{extracted equity}
+
\text{monthly cash flow}.
\end{equation}

This is the first deal as Alex presents it: not income earned once and spent, but capital placed into an asset that could later support both liquidity and continuing ownership.

\begin{figure}
\centering
\begin{tikzpicture}[
  box/.style={draw, rounded corners, align=center, text width=0.78\linewidth, minimum height=1.0cm},
  arrow/.style={->, thick}
]
\node[box] (cash) {Stated first-deal cash\\\(C_0 \approx \$28{,}000\)};
\node[box, below=0.55cm of cash] (acq) {Acquire and retain\\the real estate asset};
\node[box, below=0.55cm of acq] (value) {Later stated value\\\(V_{\text{later}} \approx \$400{,}000\)};
\node[box, below=0.55cm of value] (extract) {About \(\$150{,}000\)\\pulled out};
\node[box, below=0.55cm of extract] (keep) {Still owned\\plus monthly rent cash flow};

\draw[arrow] (cash) -- (acq);
\draw[arrow] (acq) -- (value);
\draw[arrow] (value) -- (extract);
\draw[arrow] (extract) -- (keep);
\end{tikzpicture}
\caption{Transcript-derived reconstruction of the first-deal mechanism. The exact financing structure is not specified in the interview.}
\end{figure}

\section{Partnership: Same Values, Different Work}

The host then turns from the deal to the partnership. Alex does not describe the partnership as a formal design chosen in advance. He says they started doing deals together, started working together, and learned that it worked. He was learning from Matt and doing legwork, taking action, and trying to find deals while Matt had his own commercial brokerage and property-management context.

The rule that emerges is two-sided. On one side, the partners must align morally, spiritually, and in values. On the other side, they should not be duplicates. Alex says Matt is strong at sales, networking, meeting people, relationships, ideas, and vision. Alex says he enjoys putting things into action, running the business, running operations, and growing and scaling operations.

\begin{table}
\centering
\begin{tabular}{|p{0.42\linewidth}|p{0.42\linewidth}|}
\hline
Matt Teifke & Alex Kaufman \\
\hline
Sales & Operations \\
Networking & Execution \\
Relationships & Running the business \\
Ideas and vision & Growth and scale \\
\hline
\end{tabular}
\caption{The partnership roles as described in the interview.}
\end{table}

We can state the partnership condition as a compact operating rule:

\begin{equation}
\text{strong partnership}
\approx
\text{shared values}
+
\text{trust}
+
\text{different skills}.
\end{equation}

Alex is careful about one point: being close like brothers does not automatically make people good business partners. The strength came from having both brother-like trust and business complementarity.

\section{Risk, Reinvestment, And The Cost Of Inaction}

The next question asks for Alex's best and worst financial decisions. His best decision, he says, was to put every penny he made into real estate and into the businesses. His worst decision was not taking more risks.

That answer leads directly to the local puzzle of the section: how can not taking risk be a risk?

\subsection{Question \& Answer}

\paragraph{Question.}
How can avoiding risk still be risky?

\paragraph{Answer.}
Avoiding an investment can reduce one kind of risk: the risk of losing the money committed to that investment. But it creates another exposure: the lost upside from not acting. Alex states the idea plainly: when you do not put money into an investment, you may be risking the rewards you would have received from the upside.

A cautious mathematical sketch is

\begin{equation}
\text{risk of action}
=
\text{possible downside},
\end{equation}

while

\begin{equation}
\text{risk of inaction}
=
\text{forfeited upside}.
\end{equation}

Neither term is automatically larger. The point is that inaction is not neutral. It has a payoff structure too.

Alex's own rule is more aggressive than a textbook portfolio rule. He says they put everything on the line at all times and operate with complete belief that they will make it work. The chapter should preserve that as his operating claim, not as general financial advice. The general mechanism is opportunity cost; the personal style is high conviction.

\begin{equation}
\text{reinvestment rule}
:
\quad
\text{income}
\longrightarrow
\text{real estate and business}
\longrightarrow
\text{more ownership capacity}.
\end{equation}

This is how the lecture's financial logic connects back to the first deal. Money is not treated as a stopping point. It is treated as fuel for more assets, more operations, and more deals.

\section{Vision, Sacrifice, And What The Number Is For}

When asked where he sees himself in ten years, Alex does not give a unit count, cash-flow target, revenue target, or profit target. He explicitly says they do not have that kind of number. The goal is instead a community and network of entrepreneurs under Teifke Real Estate across the world, doing deals, helping each other, and growing together. He jokes that when Elon gets to Mars, they will do it on Mars too.

This is the point where the numerical story changes register. Earlier, the numbers mattered because they showed the asset mechanism. Here, the absence of a number matters. The target is not

\begin{equation}
\text{goal} = \text{fixed unit count or fixed revenue}.
\end{equation}

It is closer to

\begin{equation}
\text{goal}
=
\text{network}
+
\text{community}
+
\text{deal flow}
+
\text{mutual help}.
\end{equation}

The final advice returns to personal responsibility. Alex says that people's situations differ. When he was working at Firehouse Subs and other jobs, he had no wife, no children, and no one depending on him. That gave him room to take actions that someone with heavier obligations might not be able to take. But his broader claim is still demanding: if one wants something badly enough, one has to figure out a way, sacrifice for it, stop letting outside voices define what is possible, and become fully accountable to oneself.

So the chapter ends where it began, with redirection. The raw desire to make money is not enough. It has to be placed inside legality, sacrifice, ownership, trust, operating skill, and responsibility.

\section{Summary}

Alex Kaufman's interview is built around a sequence of conversions. Personal disorder is converted into sobriety. Minimum-wage work is converted into the search for legal self-directed earning. Trust in Matt Teifke is converted into a first real estate deal. A stated \(\$28{,}000\) first-deal figure becomes, in Alex's telling, a property later worth close to \(\$400{,}000\), with about \(\$150{,}000\) pulled out while the asset remains owned and rented. A friendship becomes a partnership when shared values combine with different skills. Risk becomes two-sided: the possibility of loss on one side, and forfeited upside on the other.

The durable lesson for the larger book is not that every beginner should copy the exact deal. The transcript does not give enough financing detail for that. The lesson is more structural: wealth in this story came from moving from labor income toward ownership, from one-lane work toward multiple earning channels, from isolated ambition toward trusted partnership, and from passive safety toward calculated exposure to upside.